\documentclass{report}
\usepackage{amsmath,amssymb}
\setlength{\parindent}{0mm}
\setlength{\parskip}{1em}
\begin{document}
\begin{center}
\rule{6in}{1pt} \
{\large Pieter Ghysels \\
{\bf Hiding global communication in the GMRES algorithm}}

Middelheimcampus MG113 \\ Middelheimlaan 1 \\ 2020 Antwerp \\ Belgium
\\
{\tt pieter.ghysels@ua.ac.be}\\
Tom Ashby\\
Karl Meerbergen\\
Wim I. Vanroose\end{center}

In many Krylov solvers running on current large scale compute clusters,
global reductions for both inner products and norms is starting to become
a bottleneck; a problem which will only get worse as we approach the era
of exascale machines. Although flop rates keep increasing through
additional parallelism, link latencies hit their physical limits. As
compute nodes are often distributed over large datacenters, these
latencies start to play an important role in algorithm design.
Furthermore, load imbalance, system noise and processor variability (due
to extreme reduction in scale and voltage) will make global communication
an even more costly global \emph{synchronization} step~[1].

Recently there has been much interest in so called $s$-step Krylov
methods~[2] that expand the Krylov base with $s$ new vectors at once
before performing an orthogonalization step. This reduces the number of
global communications by a factor $s$.

In this work, we propose the use of non-blocking or asynchronous global
collective operations to hide the latency of global communication in the
GMRES algorithm. A standard GMRES iteration~[3], based on classical
(iterated) Gram-Schmidt orthogonalization, requires at least two global
reductions, one for orthogonalization and one for normalization. First of
all, we show that this can be reduced to just a single reduction. To
avoid stability problems, the matrix vector product $z_{i+1} = Av_i$ is
modified with a shift $\sigma_{i}$ as $z_{i+1} = \left( A - \sigma_i I
\right) v_{i}$. Proper choices of the shifts are discussed. Furthermore,
we present a class of pipelined GMRES algorithms that allow overlapping
of the inner product latency with other computations and communications.
Depending on the total reduction latency, a larger pipelining depth $l$
can be used. The resulting algorithm keeps two sets of Krylov base
vectors \[ V_{i-l+1}=[v_0,v_1,\ldots v_{i-l}] \quad \text{and} \quad
Z_{i+1}=[z_0,z_1, \ldots z_{i-l},z_{i-l+1},\ldots z_{i}] \, , \] with
$V^TV = I$. Iteration $i$ expands the $V_{i-l+1}$ base using only results
of reductions started $l$ iterations ago while the $Z_{i+1}$ base is
expanded by a vector created by the (shifted) matrix vector product. To
prevent a blow-up of the condition number of the $Z_i$ base, a
\emph{correction} is applied to $Z_{i+1}$, which introduces some
redundant calculations compared to the original GMRES algorithm.

The numerical stability of the presented GMRES variation is examined
using a collection of test matrices from applications. To assess the
scalability and overall performance of the pipelined algorithms, we
developed an analytical performance model, which allows us to study
scaling aspects and determine bottlenecks on large numbers of nodes. As
non-blocking collective operations are proposed for the upcoming MPI-3
standard, they will likely become available in most MPI libraries soon.

We compare our approach with $s$-step or \emph{communication-avoiding}
GMRES, as presented in~[2]. We present some preliminary results on
pipelining applied to the CG algorithm, for which the results look
promising since the short recurrence keeps the amount of redundant
calculations small.

In the pipelined GMRES algorithms, data dependencies between subsequent
steps in the algorithm have been relaxed. As a result, the network
hardware constraints can be relaxed. Apart from hiding global reduction
latencies on large distributed clusters, this could also allow more
efficient fine grained work scheduling on current multicore and
heterogeneous (CPU + accelerator) architectures.


\textbf{References}
\begin{itemize}
\item[1.] T.~Hoefler, T.~Schneider, A.~Lumsdaine, \emph{Characterizing
the Influence of System Noise on Large-Scale Applications by Simulation}.
International Conference for High Performance Computing, Networking,
Storage and Analysis (SC'10), Nov, 2010.

\item[2.] M.~Hoemmen, \emph{Communication-avoiding Krylov subspace
methods}. PhD, University of California, 2010.

\item[3.] Y.~Saad, M.H.~Schultz, \emph{GMRES: A generalized minimal
residual algorithm for solving nonsymmetric linear systems}. SIAM J. Sci.
Stat. Comput., 7:3, 856-869, 1986.
\end{itemize}


\end{document}
